% Blueprint for "Magic and communication complexity"
% Girish, May, Parham, Yuen, arXiv:2510.07246v1.
%
% Two-kind model:
%   * Quantum-computation and communication-complexity notions are NOT in
%     Mathlib/cslib, so they are NOVEL nodes (full statements/proofs, no
%     \mathlibok). The only library leaf used here is the binary entropy
%     function.
%   * Results cited from external references (Speelman [7], the garden-hose
%     lemmas, the external communication lower bounds) are stated precisely and
%     flagged \notready with a TODO, since the paper does not reprove them.

\chapter{Preliminaries: quantum and algebraic background}

% --- The single genuine library leaf. ----------------------------------------

\begin{definition}[Binary entropy function]
  \label{def:binary-entropy}
  \lean{Real.binEntropy}
  \mathlibok
  % TODO: confirm Lean decl name; Mathlib's Real.binEntropy uses natural log
  % (negMulLog), the paper uses base-2 log; they differ by a constant factor.
  The binary entropy function is
  \[
    h(\epsilon) = -\epsilon \log \epsilon - (1-\epsilon)\log(1-\epsilon),
    \qquad \epsilon \in [0,1],
  \]
  with the convention $0\log 0 = 0$. (Recalled only so later magic-count lower
  bounds can refer to it.)
\end{definition}

% --- Novel quantum-computation primitives (not in the libraries). -------------

\begin{definition}[Pauli operators and Pauli strings]
  \label{def:pauli}
  The single-qubit Pauli operators are
  \[
    X = \begin{pmatrix} 0 & 1 \\ 1 & 0 \end{pmatrix}, \qquad
    Z = \begin{pmatrix} 1 & 0 \\ 0 & -1 \end{pmatrix}, \qquad
    Y = iXZ.
  \]
  They satisfy $X^2 = Z^2 = I$ and the anticommutation relation $ZX = -XZ$.
  For $\vec a = (a_1,\dots,a_n), \vec b \in \{0,1\}^n$ an ($n$-qubit) \emph{Pauli
  string} is
  \[
    X^{\vec a} Z^{\vec b}
      = \bigotimes_{i=1}^n X^{a_i} Z^{b_i},
  \]
  possibly multiplied by a global phase in $\{\pm 1, \pm i\}$. The set of all
  such operators (with phases) forms the $n$-qubit Pauli group. We freely use
  $X^{\vec x}\ket{0\cdots 0} = \ket{\vec x}$ for $\vec x\in\{0,1\}^n$.
\end{definition}

\begin{definition}[Phase gate]
  \label{def:phase-gate}
  The phase gate is $P = \begin{pmatrix} 1 & 0 \\ 0 & i \end{pmatrix}$. It
  satisfies the conjugation relation $P X P^\dagger = i X Z$ (so conjugating a
  Pauli $X$ by $P$ produces a Pauli up to phase) and $PZP^\dagger = Z$.
\end{definition}

\begin{definition}[Clifford group and Clifford circuit]
  \label{def:clifford}
  \uses{def:pauli}
  A unitary $C$ on $n$ qubits is a \emph{Clifford} operator if it normalizes the
  Pauli group, i.e.\ $C\,Q\,C^\dagger$ is a Pauli string (up to global phase) for
  every Pauli string $Q$. The Clifford group is generated by the Hadamard gate
  $H$, the phase gate $P$ (Definition~\ref{def:phase-gate}), and $\mathrm{CNOT}$.
  A \emph{Clifford circuit} is a circuit composed only of Clifford gates.
\end{definition}

\begin{definition}[Magic gate, weight, $T$ gate, Clifford+$T$]
  \label{def:magic-gate}
  \uses{def:clifford}
  A \emph{magic gate} (equivalently, non-Clifford gate) is any unitary that is
  not a Clifford operator. The \emph{weight} of a gate is the number of qubits on
  which it acts. The \emph{$T$ gate} is
  $T = \begin{pmatrix} 1 & 0 \\ 0 & e^{i\pi/4}\end{pmatrix}$, a weight-$1$ magic
  gate; the \emph{Clifford$+T$} gate set is $\{H,P,\mathrm{CNOT},T\}$. A
  \emph{Clifford$+$Magic circuit} is a circuit composed of Clifford gates and
  magic gates. The $T$ gate satisfies $T X T^\dagger = e^{-i\pi/4} P X$ and
  $TZ = ZT$ (it commutes with Pauli $Z$).
\end{definition}

\begin{definition}[$T$-depth and magic-depth]
  \label{def:tdepth}
  \uses{def:magic-gate}
  Write a Clifford$+T$ circuit as a sequence of \emph{layers}, where each layer
  is either a Clifford circuit or a layer of (parallel) $T$ gates. The
  \emph{$T$-depth} of the circuit is the number of $T$-layers in such a
  decomposition; the \emph{$T$-depth of a unitary} $U$ is the minimum $T$-depth
  over Clifford$+T$ circuits implementing $U$. The \emph{magic-depth} is defined
  analogously with general magic-gate layers in place of $T$-layers.
\end{definition}

\begin{definition}[Density operator, trace norm, trace distance]
  \label{def:trace-norm}
  A \emph{density operator} $\rho$ is a positive semidefinite operator with
  $\Tr\rho = 1$. The \emph{trace norm} of an operator $A$ is
  $\onenorm{A} = \Tr\!\sqrt{A^\dagger A}$, and the \emph{trace distance} between
  density operators $\rho,\sigma$ is $\tfrac12\onenorm{\rho-\sigma}$.
\end{definition}

\begin{definition}[Diamond norm of a superoperator]
  \label{def:diamond-norm}
  \uses{def:trace-norm}
  For a Hermicity-preserving superoperator $\Phi$ (e.g.\ a difference of quantum
  channels), the \emph{diamond norm} is
  \[
    \dnorm{\Phi}
      = \max_{\rho} \onenorm{(\Phi \otimes \mathrm{id})(\rho)},
  \]
  the maximum taken over density operators $\rho$ on the input space tensored
  with an ancilla of the same dimension. For unitaries $U,U'$ we write
  $\dnorm{U-U'}$ for the diamond norm of the difference of the channels
  $\rho \mapsto U\rho U^\dagger$ and $\rho\mapsto U'\rho U'^\dagger$.
\end{definition}

\begin{definition}[Parity function]
  \label{def:parity}
  \uses{def:pauli}
  For subsets $S_A, S_B \subseteq [n]$, the associated \emph{parity function} on
  inputs $x,y \in \{0,1\}^n$ is
  \[
    p(x,y) = \sum_{i \in S_A} x_i + \sum_{i \in S_B} y_i \pmod 2 .
  \]
  More generally a parity function of $(x,y)$ is any $\mathbb{F}_2$-affine
  function of the bits of $(x,y)$.
\end{definition}

\begin{definition}[EPR pairs, Bell-basis measurement, teleportation]
  \label{def:teleportation}
  \uses{def:pauli}
  An \emph{EPR pair} is the two-qubit state
  $\tfrac{1}{\sqrt2}(\ket{00}+\ket{11})$. A \emph{Bell-basis measurement} on two
  qubits projects onto the four maximally entangled states
  $(X^a Z^b \otimes I)\tfrac{1}{\sqrt2}(\ket{00}+\ket{11})$, $a,b\in\{0,1\}$,
  returning the outcome $(a,b)$. \emph{Teleportation}: if Alice holds a state
  $\ket{\phi}$ and one half of an EPR pair whose other half is held by Bob, then
  a Bell-basis measurement by Alice on $\ket{\phi}$ and her EPR half, with
  outcome $(a,b)$, leaves Bob's half in the state $X^a Z^b \ket{\phi}$. The
  outcome $(a,b)$ is uniformly random and independent of $\ket{\phi}$.
\end{definition}

\chapter{Communication models}

\begin{definition}[Simultaneous message passing protocol; the $\parallel$ models]
  \label{def:smp}
  Let $f : \{0,1\}^n \times \{0,1\}^n \to \{0,1\}$ be a (partial or total)
  Boolean function and $\epsilon \in [0,1]$. A \emph{simultaneous message
  passing} protocol $P$ for $f$ has three parties: Alice (input $x$), Bob (input
  $y$), and a referee. Alice and Bob send the referee message systems $M_A, M_B$
  respectively, the referee outputs a bit $c = P(x,y)$, and the parties do not
  otherwise communicate. The protocol is \emph{$\epsilon$-correct} if
  $\Pr[P(x,y)=f(x,y)] \ge 1-\epsilon$ for all $(x,y)$ in the support of $f$. The
  \emph{cost} $\mathrm{cost}(P)$ is the total number of bits (classical) or qubits
  (quantum) sent by Alice and Bob; the complexity of $f$ is the minimum cost over
  $\epsilon$-correct protocols. The variants are:
  \begin{itemize}
    \item messages classical = $\Rpar$, quantum = $\Qpar$;
    \item $\Dpar$ = classical, deterministic, $\epsilon = 0$;
    \item shared \emph{public} randomness $r$ available to all three parties:
      superscript $\mathsf{pub}$, giving $\Rparpub_\epsilon, \Qparpub_\epsilon$;
    \item Alice and Bob share \emph{entanglement}: superscript $*$, giving
      $\Rpars, \Qpars$.
  \end{itemize}
  When the subscript $\epsilon$ is dropped, $\epsilon = 1/3$. The complexities are
  $\Dpar(f), \Rpar_\epsilon(f), \Qpar_\epsilon(f), \Rparpub_\epsilon(f),
  \Qparpub_\epsilon(f), \Rpars(f), \Qpars(f)$.
\end{definition}

\begin{definition}[Private simultaneous message (PSM) task]
  \label{def:psm}
  \uses{def:trace-norm, def:smp}
  A \emph{private simultaneous message} task is given by a Boolean function
  $f:\{0,1\}^n\times\{0,1\}^n\to\{0,1\}$ and parameters $\epsilon,\delta\in[0,1]$.
  Alice (input $x$) sends a message system $M_0$ and Bob (input $y$) sends $M_1$
  to the referee; from $M = M_0 M_1$ the referee prepares an output bit $z$ on a
  system $Z$. Writing $\rho_M(x,y)$ for the density operator on $M$ produced on
  inputs $(x,y)$ and $f_{x,y} = f(x,y)$, the task requires:
  \begin{itemize}
    \item \emph{$\epsilon$-correctness}: there is a decoding map
      $\mathbf{V}_{M\to Z\bar M}$ with, for all $(x,y)$ in the support of $f$,
      \[
        \onenorm{\Tr_{\bar M}\!\big(\mathbf{V}_{M\to Z\bar M}\,\rho_M(x,y)\,
          \mathbf{V}^\dagger_{M\to Z\bar M}\big) - \ketbra{f_{x,y}}{f_{x,y}}_Z}
          \le \epsilon ;
      \]
    \item \emph{$\delta$-security}: there is a simulator channel
      $\mathcal{S}_{Z\to M}$ with, for all $(x,y)$ on which $f$ is defined,
      \[
        \onenorm{\rho_M(x,y) - \mathcal{S}_{Z\to M}
          (\ketbra{f_{x,y}}{f_{x,y}}_Z)} \le \delta .
      \]
  \end{itemize}
  Classical messages give $\PSM$; quantum messages give $\PSQM$; shared
  entanglement adds the superscript $*$ giving $\PSMs, \PSQMs$. The cost is the
  total number of bits (resp.\ qubits) sent; $\PSM_{\epsilon,\delta}(f)$ etc.\
  denote the minimal cost over $\epsilon$-correct, $\delta$-secure protocols (with
  $\epsilon=1/3$ when omitted).
\end{definition}

\begin{definition}[Non-local quantum computation (NLQC)]
  \label{def:nlqc}
  \uses{def:teleportation}
  In the \emph{non-local quantum computation} setting there is no referee:
  Alice and Bob each send a single simultaneous quantum message to the other and
  then perform local operations. Concretely, to implement a target channel
  $\mathcal{N}$ on a bipartite input, Alice and Bob apply local channels
  $\mathcal{V}^L,\mathcal{V}^R$ to their inputs together with their halves of
  shared entanglement, exchange messages, and apply further local channels
  $\mathcal{W}^L,\mathcal{W}^R$, so that the overall map equals $\mathcal{N}$.
  The cost is the total communication and shared entanglement.
\end{definition}

\begin{definition}[Two-way classical communication complexity]
  \label{def:two-way}
  Let $f : \{0,1\}^n \times \{0,1\}^n \to \{0,1\}$ and $\epsilon \in [0,1]$. A
  \emph{two-way classical communication} protocol involves Alice (input $x$) and
  Bob (input $y$), who may share a random string $r$. They exchange a sequence of
  messages (each computed from the local input, $r$, and previously received
  messages), and Alice finally outputs a bit $z$. The protocol is
  $\epsilon$-correct if $\Pr[f(x,y)=z]\ge 1-\epsilon$ for all $(x,y)$ on which $f$
  is defined. The cost is the number of bits exchanged, maximized over inputs;
  the \emph{two-way complexity} $\Rtwo_\epsilon(f)$ is the minimum cost over
  $\epsilon$-correct protocols.
\end{definition}

\chapter{Magic lower bounds from communication complexity}

\section{Computation and communication models}

\begin{definition}[Unitary Clifford+Magic circuit and $\Mu$]
  \label{def:unitary-cm}
  \uses{def:magic-gate}
  A \emph{unitary Clifford$+$Magic} circuit is a circuit of Clifford gates and
  magic gates. We say it computes $f:\{0,1\}^n\to\{0,1\}$ with correctness
  $\epsilon$ if there is an advice state $\ket{\psi}$ (independent of the input)
  such that running the circuit on $\ket{x}\ket{\psi}$ and measuring the first
  qubit returns $f(x)$ with probability $\ge 1-\epsilon$ for all $x$ in the
  support of $f$. The cost $\Mu_{\epsilon,c_M}(f)$ is the minimal number of magic
  gates, each of weight at most $c_M$, in any such circuit. Access to an arbitrary
  (arbitrarily large) advice state is allowed.
\end{definition}

\begin{definition}[Mixed Clifford+Magic circuit and $\Mmix$]
  \label{def:mixed-cm}
  \uses{def:unitary-cm}
  A \emph{mixed Clifford$+$Magic} circuit is a quantum operation
  \[
    \mathcal{N}(\cdot) = \sum_i p_i\, U_i(\cdot)U_i^\dagger ,
  \]
  where $\{p_i\}$ is a probability distribution and each $U_i$ is a unitary
  Clifford$+$Magic circuit. Its magic-gate count is the worst-case count among the
  $U_i$. It computes $f$ with correctness $\epsilon$ if (with an arbitrary advice
  state) measuring the first qubit yields $f(x)$ with probability $\ge 1-\epsilon$.
  The cost $\Mmix_{\epsilon,c_M}(f)$ is the minimal number of magic gates (of
  weight $\le c_M$) in any such operation computing $f$.
\end{definition}

\begin{definition}[Adaptive Clifford+Magic circuit and $\Mad$]
  \label{def:adaptive-cm}
  \uses{def:unitary-cm}
  An \emph{adaptive Clifford$+$Magic} circuit is composed of Clifford gates,
  arbitrary magic gates, and mid-circuit computational-basis measurements, where
  later gate choices may depend arbitrarily on earlier measurement outcomes. Its
  cost is the total number of magic gates plus measurements in the worst-case
  run. The cost $\Mad_{\epsilon,c_M}(f)$ is the minimal such cost over adaptive
  circuits using magic gates and measurements each acting on at most $c_M$ qubits
  (a single computational-basis measurement of $c_M$ qubits counts once) that
  compute $f$ with probability $\ge 1-\epsilon$, with an arbitrary advice state.
\end{definition}

\begin{lemma}[Comparison of computational models]
  \label{lem:model-comparison}
  \uses{def:unitary-cm, def:mixed-cm, def:adaptive-cm}
  For every Boolean function $f$, $\Mu_{\epsilon}(f) \ge \Mmix_{\epsilon}(f)$ and
  $\Mu_{\epsilon}(f) \ge \Mad_{\epsilon}(f)$. Moreover the adaptive model can
  simulate the mixed model (by preparing $\ket{+}$ states and measuring), though
  the costs may be incomparable because $\Mad$ also counts measurements.
\end{lemma}
\begin{proof}
  \uses{def:unitary-cm, def:mixed-cm, def:adaptive-cm}
  A unitary Clifford$+$Magic circuit is the special case of a mixed circuit with
  a single term ($p_1 = 1$), and a special case of an adaptive circuit using no
  measurements; in both cases the magic-gate count is unchanged. Hence any
  unitary circuit witnessing $\Mu_\epsilon(f)$ is admissible in the mixed and
  adaptive models, giving $\Mu_\epsilon(f)\ge\Mmix_\epsilon(f)$ and
  $\Mu_\epsilon(f)\ge\Mad_\epsilon(f)$. For the simulation claim, an adaptive
  circuit can prepare $\ket{+}$ states and measure in the computational basis to
  generate randomness and then condition subsequent gates on the outcomes,
  reproducing the random choice $i$ with distribution $\{p_i\}$ of a mixed
  circuit; but $\Mad$ additionally charges for these measurements, so the two
  costs need not be comparable.
\end{proof}

\section{Lower bounds on magic gate count from communication complexity}

\begin{lemma}[A single parity has $\Dpar$ cost $2$]
  \label{lem:parity-dpar}
  \uses{def:parity, def:smp}
  For any parity function $p(x,y) = \sum_{i\in S_A} x_i + \sum_{i\in S_B} y_i
  \pmod 2$, the referee can compute $p(x,y)$ in the $\Dpar$ model with
  communication cost $2$.
\end{lemma}
\begin{proof}
  \uses{def:parity, def:smp}
  Alice sends the single bit $\sum_{i\in S_A} x_i \bmod 2$ and Bob sends the
  single bit $\sum_{i\in S_B} y_i \bmod 2$. The referee XORs the two received
  bits to obtain $p(x,y)$. Two bits are communicated.
\end{proof}

\begin{lemma}[Pauli corrections are parities of the input]
  \label{lem:pauli-parity}
  \uses{def:clifford, def:pauli, def:parity}
  Let $C$ be a Clifford circuit on $n$ qubits and let $X^{\vec a}Z^{\vec b}$ be a
  Pauli string whose exponents $\vec a,\vec b$ are $\mathbb{F}_2$-affine functions
  (parity functions) of an input $(x,y)$. Then $C\,(X^{\vec a}Z^{\vec b})\,
  C^\dagger$ is a Pauli string (up to phase) each of whose $X$- and $Z$-exponents
  is again a parity function of $(x,y)$, with the parities determined by $C$ and
  not by the input.
\end{lemma}
\begin{proof}
  \uses{def:clifford, def:pauli}
  Conjugation by a Clifford $C$ maps each Pauli generator $X_i, Z_i$ to a Pauli
  string (Definition~\ref{def:clifford}); on exponent vectors over
  $\mathbb{F}_2^{2n}$ this action is $\mathbb{F}_2$-linear (the symplectic
  representation of the Clifford group), depending only on $C$ and ignoring
  global phases. Hence $C(X^{\vec a}Z^{\vec b})C^\dagger$ has exponent vector
  $L_C(\vec a, \vec b)$ for a fixed linear map $L_C$ determined by $C$. If
  $(\vec a,\vec b)$ are affine functions of $(x,y)$, then so is
  $L_C(\vec a,\vec b)$, i.e.\ every output exponent is a parity function of
  $(x,y)$ with coefficients fixed by $C$.
\end{proof}

\begin{theorem}[Theorem 7: $\Dpar$ vs unitary magic count]
  \label{thm:dpar-unitary}
  \uses{def:unitary-cm, def:smp}
  Let $f$ be a Boolean function. For all $\epsilon < \tfrac12$ and all $c_M$,
  \[
    \frac{1}{4 c_M}\big(\Dpar(f) - 2\big) \le \Mu_{\epsilon}(f),
  \]
  where the magic gates have weight at most $c_M$.
\end{theorem}
\begin{proof}
  \uses{lem:pauli-parity, lem:parity-dpar, def:unitary-cm, def:smp}
  Take a unitary Clifford$+$Magic circuit computing $f$ with probability
  $p = 1-\epsilon > 1/2$, with $k$ magic gates each of weight $\le c_M$.
  Decompose it into layers $C^0, M_1, C^1, M_2, \dots, M_k, C^k$ where each $C^j$
  is a Clifford circuit and each $M_j$ is a magic gate. Run on input
  $\ket{z}\ket{\psi}$ with $z$ split into $(x,y)$; view the input as
  $X_A^{\vec x}\ket{x}_A X_B^{\vec y}\ket{y}_B\ket{\psi}_E$ obtained from the
  all-zeros state.

  Give a $\Dpar$ protocol. The referee runs the circuit on
  $\ket{0}\ket{0}\ket{\psi}$. After the first Clifford layer $C^0$, the conjugated
  input Paulis become a Pauli correction $\sigma_{ABE}[x,y]$ on the state, and by
  Lemma~\ref{lem:pauli-parity} each of its $X$- and $Z$-exponents is a parity
  function of $(x,y)$. Before applying $M_1$ we must undo the Pauli corrections on
  the (at most $c_M$) wires that $M_1$ acts on; each such wire carries a
  correction $X^a Z^b$ with $a,b$ parity functions of $(x,y)$, i.e.\ at most
  $2c_M$ parity functions. By Lemma~\ref{lem:parity-dpar} each parity costs $2$
  bits of $\Dpar$ communication. Apply $M_1$, then the next Clifford layer; the
  remaining corrections conjugate through to new parity-valued corrections
  (Lemma~\ref{lem:pauli-parity}), and we repeat before each $M_j$. Finally one
  more parity determines whether there is an $X$ correction on the output qubit
  before the measurement (a $Z$ correction does not affect the computational-basis
  measurement and may be left uncorrected). Correcting and measuring yields a
  sample from the circuit's output distribution.

  Repeating with the same parity values, the referee estimates the output
  distribution; since the circuit is correct with probability $1-\epsilon > 1/2$,
  outputting the majority bit is correct. The total cost is at most $4c_M$ per
  magic gate (at most $2c_M$ parities, $2$ bits each) plus $2$ for the final
  parity, so $\Dpar(f) \le 4c_M\cdot \Mu_{\epsilon<1/2}(f) + 2$. Rearranging gives
  the claim.
\end{proof}

\begin{proposition}[Remark 8: post-selection strengthening]
  \label{prop:postselection}
  \uses{thm:dpar-unitary}
  $\Dpar(f)$ also lower bounds the number of magic gates plus the number of
  post-selected qubits in any Clifford$+$Magic$+$post-selection circuit computing
  $f$: precisely, the simulation of Theorem~\ref{thm:dpar-unitary} extends so
  that each qubit of post-selection adds $2$ bits of $\Dpar$ communication.
\end{proposition}
\begin{proof}
  \uses{thm:dpar-unitary, lem:parity-dpar}
  Run the simulation of Theorem~\ref{thm:dpar-unitary}. Just before a
  post-selection onto a fixed state, send the (at most $2$ per qubit) parity-valued
  Pauli corrections acting on the post-selected qubit, exactly as is done before a
  magic gate; by Lemma~\ref{lem:parity-dpar} this costs $2$ bits per qubit of
  post-selection. The referee corrects and then post-selects. Thus $\Dpar(f)$ also
  lower bounds the magic-gate count plus the number of post-selected qubits.
\end{proof}

\begin{theorem}[Theorem 9: $\Rparpub$ vs mixed magic count]
  \label{thm:rpar-mixed}
  \uses{def:mixed-cm, def:smp}
  Let $f$ be a Boolean function. For all $\epsilon$ and $c_M$,
  \[
    \frac{1}{4 c_M}\big(\Rparpub_\epsilon(f) - 2\big) \le \Mmix_{\epsilon}(f).
  \]
\end{theorem}
\begin{proof}
  \uses{thm:dpar-unitary, def:mixed-cm, def:smp, lem:parity-dpar}
  Take a mixed Clifford$+$Magic circuit $\mathcal{N}(\cdot)=\sum_i p_i U_i(\cdot)
  U_i^\dagger$ computing $f$ with probability $\ge 1-\epsilon$; let $P_i(z)$ be the
  probability that $U_i$ outputs $f(z)$ on input $z$, so the success probability
  is $p_{\mathrm{suc}}(z) = \sum_i p_i P_i(z) \ge 1-\epsilon$. The referee, Alice
  and Bob use public randomness to sample $i$ with probability $p_i$. They then
  run the $\Dpar$ protocol of Theorem~\ref{thm:dpar-unitary} for the unitary
  $U_i$, but the referee samples the final measurement \emph{once} and outputs
  the result rather than taking a majority. The protocol is then correct with
  probability $P_i$, and overall with probability $\sum_i p_i P_i \ge 1-\epsilon$.
  The communication for the selected $U_i$ is at most $4c_M$ times the number of
  magic gates in $U_i$ plus $2$ (Theorem~\ref{thm:dpar-unitary},
  Lemma~\ref{lem:parity-dpar}); the worst case over $i$ is $4c_M
  \Mmix_{\epsilon,c_M}(f) + 2$. Hence $\Rparpub_\epsilon(f) \le
  4c_M\Mmix_{\epsilon}(f)+2$, which rearranges to the claim.
\end{proof}

\begin{theorem}[Theorem 10: two-way $\Rtwo$ vs adaptive magic count]
  \label{thm:r-adaptive}
  \uses{def:adaptive-cm, def:two-way}
  Let $f$ be a Boolean function. For all $\epsilon$ and $c_M$,
  \[
    \frac{1}{2 c_M}\big(\Rtwo_\epsilon(f) - 1\big) \le \Mad_{\epsilon,c_M}(f).
  \]
\end{theorem}
\begin{proof}
  \uses{lem:pauli-parity, def:adaptive-cm, def:two-way}
  Build a two-way protocol from an adaptive circuit. Decompose the circuit into
  layers, each an arbitrary Clifford circuit followed by at most one magic gate or
  one (computational-basis) measurement, occurring first in the layer. Alice
  prepares the advice state $\ket{\psi}_E$ and runs on
  $\ket{x}_A X_B^{\vec y}\ket{y}_B\ket{\psi}_E$; she applies the first Clifford
  layer, producing a Pauli correction $\sigma_{ABE}[y]$ whose exponents are
  parities of the input (Lemma~\ref{lem:pauli-parity}).

  If the first non-Clifford operation is a magic gate (weight $\le c_M$), Bob
  computes the identities of the Pauli corrections on the $\le c_M$ wires it acts
  on and sends them to Alice: for each wire there may be an $X$ and a $Z$
  correction, so $2c_M$ bits. Alice corrects and applies the gate. If instead it
  is a measurement (of $\le c_M$ qubits), Bob sends the $X$ corrections on the
  measured wires, Alice corrects, measures, and sends the $\le c_M$ outcome bits
  back to Bob; total $2c_M$ bits. In either case both players learn the
  measurement outcomes and choose the next layer accordingly. Repeating for every
  layer costs $2c_M$ per magic gate or measurement. A final Pauli-$X$ correction
  before the output measurement contributes $+1$. The outcome equals $f(x,y)$ with
  probability $1-\epsilon$, so $\Rtwo_\epsilon(f) \le 2c_M
  \Mad_{\epsilon,c_M}(f)+1$, giving the claim.
\end{proof}

\subsection*{Bounds from parity decision trees}

\begin{definition}[Definition 11: non-adaptive parity decision tree]
  \label{def:pdt}
  \uses{def:parity}
  A \emph{non-adaptive parity decision tree} ($\PDTna$) of depth $k$ computing
  $f:\{0,1\}^n\to\{0,1\}$ is a function $g:\{0,1\}^k\to\{0,1\}$ together with
  parity functions $p_1,\dots,p_k$ such that $f(x) = g(p_1(x),\dots,p_k(x))$ for
  all $x$. The $\PDTna(f)$ complexity is the minimal such $k$.
\end{definition}

\begin{definition}[Definition 12: randomized non-adaptive parity decision tree]
  \label{def:rpdt}
  \uses{def:pdt}
  A \emph{randomized non-adaptive parity decision tree} ($\RPDTna$) of depth $k$
  computing $f$ with error $\epsilon$ is a probability distribution over
  $\PDTna$'s each of depth $\le k$ such that, for every input $x$, the sampled
  tree outputs $f(x)$ with probability $\ge 1-\epsilon$. $\RPDTna_\epsilon(f)$ is
  the minimal such $k$.
\end{definition}

\begin{proposition}[Eq.\ (18): $\PDTna$ lower bounds unitary magic count]
  \label{prop:pdt-unitary}
  \uses{def:pdt, def:unitary-cm}
  For all $c_M$, $\ \dfrac{1}{2c_M}\big(\PDTna(f) - 1\big) \le
  \Mu_{\epsilon<1/2}(f)$.
\end{proposition}
\begin{proof}
  \uses{thm:dpar-unitary, lem:pauli-parity, def:pdt}
  In the proof of Theorem~\ref{thm:dpar-unitary} the communication from Alice and
  Bob consists entirely of values of parity functions of $(x,y)$
  (Lemma~\ref{lem:pauli-parity}), and these parity functions depend only on the
  circuit, not on the input. Hence the referee computes the output as a fixed
  function of at most $2c_M$ parities per magic gate plus one final parity, i.e.\
  a non-adaptive parity decision tree of depth $2c_M \Mu_{\epsilon<1/2}(f) + 1$.
  Therefore $\PDTna(f) \le 2c_M\Mu_{\epsilon<1/2}(f) + 1$, which rearranges to the
  claim.
\end{proof}

\begin{proposition}[Eq.\ (19): $\PDTna$ vs unitary $T$-count]
  \label{prop:pdt-t}
  \uses{def:pdt, def:unitary-cm, def:magic-gate}
  $\ \PDTna(f) - 1 \le \Tcu_{\epsilon<1/2}(f)$, where $\Tcu$ counts $T$ gates.
\end{proposition}
\begin{proof}
  \uses{prop:pdt-unitary, def:magic-gate}
  Specialize Proposition~\ref{prop:pdt-unitary} to $T$ gates, where $c_M = 1$.
  Since $T$ commutes with Pauli $Z$ (Definition~\ref{def:magic-gate}), the $Z$
  corrections need not be undone before a $T$ gate, halving the number of required
  parities to one per $T$ gate. Hence $\PDTna(f) \le \Tcu_{\epsilon<1/2}(f) + 1$.
\end{proof}

\begin{proposition}[Eq.\ (20): $\RPDTna$ lower bounds mixed magic count]
  \label{prop:rpdt-mixed}
  \uses{def:rpdt, def:mixed-cm}
  For all $c_M$, $\ \dfrac{1}{2c_M}\big(\RPDTna_\epsilon(f) - 1\big) \le
  \Mmix_\epsilon(f)$.
\end{proposition}
\begin{proof}
  \uses{thm:rpar-mixed, prop:pdt-unitary, def:rpdt}
  For each $U_i$ in the mixed circuit, the $\Dpar$ protocol of
  Theorem~\ref{thm:dpar-unitary} defines a (deterministic) non-adaptive parity
  decision tree as in Proposition~\ref{prop:pdt-unitary}, except that the leaf now
  \emph{samples} the output from the final measurement distribution rather than
  taking a majority; this is a $\RPDTna$ that outputs $f(x)$ with the same
  probability as running $U_i$. Adding the public random choice of $i$ (with
  probability $p_i$) gives a single $\RPDTna$ outputting $f(x)$ with the same
  probability as $\sum_i p_i U_i(\cdot)U_i^\dagger$. The number of parities needed
  for the worst-case $U_i$ is $2c_M\Mmix_\epsilon(f) + 1$, so
  $\RPDTna_\epsilon(f) \le 2c_M\Mmix_\epsilon(f) + 1$.
\end{proof}

\begin{proposition}[Eq.\ (21): $\RPDTna$ vs mixed $T$-count]
  \label{prop:rpdt-t}
  \uses{def:rpdt, def:mixed-cm, def:magic-gate}
  $\ \RPDTna_\epsilon(f) - 1 \le \Tcmix_\epsilon(f)$.
\end{proposition}
\begin{proof}
  \uses{prop:rpdt-mixed, prop:pdt-t}
  Specialize Proposition~\ref{prop:rpdt-mixed} to $T$ gates ($c_M=1$) and use, as
  in Proposition~\ref{prop:pdt-t}, that $T$ commutes with $Z$ so only one parity
  per $T$ gate is needed. This gives $\RPDTna_\epsilon(f) \le \Tcmix_\epsilon(f) +
  1$.
\end{proof}

\section{Lower bounds for concrete unitary operators}

\begin{definition}[Magic count of a unitary operator]
  \label{def:magic-operator}
  \uses{def:unitary-cm, def:mixed-cm, def:diamond-norm}
  Extend $\Mu_\epsilon$ and $\Mmix_\epsilon$ to a target unitary $U$ by requiring
  the circuit to implement $U$ to within $\epsilon$ in diamond norm,
  $\dnorm{U-U'}\le\epsilon$, where $U'$ is the implemented unitary. The cost is
  the minimal magic count of such an approximating circuit.
\end{definition}

\begin{lemma}[A unitary computing $f$ inherits its magic lower bound]
  \label{lem:unitary-computes-f}
  \uses{def:magic-operator}
  If a circuit using a unitary $U$ computes $f$ with probability $1$ and no extra
  magic gates, then $\Mu_\epsilon(U) \ge \Mu_\epsilon(f)$ and $\Mmix_\epsilon(U)
  \ge \Mmix_\epsilon(f)$.
\end{lemma}
\begin{proof}
  \uses{def:magic-operator}
  Suppose $U'$ approximates $U$ with $\dnorm{U-U'}\le\epsilon$. Replacing $U$ by
  $U'$ in the circuit that computes $f$ exactly yields a circuit computing $f$
  with probability $\ge 1-\epsilon$ (the diamond-norm bound limits the change in
  any measurement outcome by $\epsilon$). Since the surrounding gates add no
  magic, any circuit implementing $U$ to error $\epsilon$ gives a circuit
  computing $f$ to error $\epsilon$ with the same magic count, so the magic count
  of $U$ is at least that of $f$ in both the unitary and mixed models.
\end{proof}

\begin{definition}[Generalized Toffoli gate]
  \label{def:toffoli}
  The \emph{$n$-qubit generalized Toffoli} gate acts by
  \[
    \Toffoli_n : \ket{x_1,\dots,x_n,b} \mapsto
      \Big\lvert x_1,\dots,x_n,\, b \oplus \textstyle\bigwedge_{i=1}^n x_i
      \Big\rangle,
  \]
  i.e.\ it XORs the AND of the first $n$ bits into the target qubit.
\end{definition}

\begin{definition}[Equality function]
  \label{def:equality}
  The \emph{equality function} $\Equal_n : \{0,1\}^n\times\{0,1\}^n\to\{0,1\}$ is
  $\Equal_n(x,y) = 1$ iff $x = y$.
\end{definition}

\begin{construction}[Computing equality from $\Toffoli_n$]
  \label{constr:toffoli-equality}
  \uses{def:toffoli, def:equality}
  On $2n+1$ qubits $A_1,\dots,A_n, B_1,\dots,B_n, C$ with input $\ket{x,y,0}$:
  \begin{enumerate}
    \item for each $i$, apply CNOT with control $A_i$, target $B_i$, giving
      $\ket{x_i, x_i\oplus y_i}$, then apply $X$ to $B_i$ giving
      $\ket{x_i, x_i\oplus y_i\oplus 1}$;
    \item apply $\Toffoli_n$ controlled on $B_1,\dots,B_n$ with target $C$.
  \end{enumerate}
  The target $C$ flips iff all $B_i = 1$, i.e.\ iff $x_i = y_i$ for all $i$; so
  the circuit computes $\Equal_n(x,y)$ using one $\Toffoli_n$ and otherwise only
  Clifford gates.
\end{construction}

\begin{lemma}[Equality $\Dpar$ lower bound]
  \label{lem:equality-dpar}
  \uses{def:equality, def:smp}
  \notready
  $\Dpar(\Equal_n) = \Omega(n)$.
\end{lemma}
% TODO: cited from Kushilevitz--Nisan, "Communication Complexity" [29]; not
% reproved in the paper.

\begin{lemma}[Equality $\Rparpub$ lower bound]
  \label{lem:equality-rpar}
  \uses{def:equality, def:smp}
  \notready
  $\Rparpub_\epsilon(\Equal_n) = \Omega(\min\{\log(1/\epsilon), n\})$.
\end{lemma}
% TODO: cited from Jacobs--Jeang--Podolskii--Prior--Volkovich [28]; not reproved.

\begin{lemma}[Lemma 13: magic lower bounds for $\Toffoli_n$]
  \label{lem:toffoli-bound}
  \uses{def:toffoli, def:magic-operator, constr:toffoli-equality}
  We have $\Mmix_\epsilon(\Toffoli_n) = \Omega(\min\{\log(1/\epsilon), n\})$ and,
  for all $\epsilon<\tfrac12$, $\Mu_\epsilon(\Toffoli_n) = \Omega(n)$.
\end{lemma}
\begin{proof}
  \uses{lem:unitary-computes-f, constr:toffoli-equality, thm:dpar-unitary,
    thm:rpar-mixed, lem:equality-dpar, lem:equality-rpar}
  By Construction~\ref{constr:toffoli-equality}, $\Toffoli_n$ computes $\Equal_n$
  with no magic overhead, so by Lemma~\ref{lem:unitary-computes-f}
  $\Mmix_\epsilon(\Toffoli_n) \ge \Mmix_\epsilon(\Equal_n)$ and
  $\Mu_\epsilon(\Toffoli_n)\ge\Mu_\epsilon(\Equal_n)$. By
  Theorem~\ref{thm:rpar-mixed},
  $\Mmix_\epsilon(\Equal_n) \ge \tfrac{1}{4c_M}(\Rparpub_\epsilon(\Equal_n)-1)$,
  and $\Rparpub_\epsilon(\Equal_n) = \Omega(\min\{\log(1/\epsilon),n\})$
  (Lemma~\ref{lem:equality-rpar}) with $c_M = O(1)$, giving the mixed bound. By
  Theorem~\ref{thm:dpar-unitary},
  $\Mu_\epsilon(\Equal_n) \ge \tfrac{1}{4c_M}(\Dpar(\Equal_n)-2)$, and
  $\Dpar(\Equal_n) = \Omega(n)$ (Lemma~\ref{lem:equality-dpar}), giving the
  unitary bound $\Omega(n)$.
\end{proof}

\begin{definition}[Quantum multiplexer]
  \label{def:multiplexer}
  The \emph{$n$-qubit quantum multiplexer} acts, for $x\in\{0,1\}^n$, an
  $\lceil\log_2 n\rceil$-bit index $i$, and $b\in\{0,1\}$, by
  \[
    \Multiplex_n : \ket{i, x, b} \mapsto
      \ket{i,\, x_1,\dots,x_{i-1}, b, x_{i+1},\dots,x_n,\, x_i} ,
  \]
  i.e.\ controlled on $i$ it swaps the $i$-th bit of $x$ with the target $b$.
\end{definition}

\begin{definition}[Index function]
  \label{def:index}
  The \emph{index function} is $\Index_n(i,x) = x_i$, where $x\in\{0,1\}^n$ and
  $i\in\{0,1\}^{\lceil\log_2 n\rceil}$.
\end{definition}

\begin{construction}[The multiplexer computes the index function]
  \label{constr:multiplexer-index}
  \uses{def:multiplexer, def:index}
  Setting the target $b = 0$, $\Multiplex_n$ swaps $x_i$ into the target
  register, so reading the target gives $\Index_n(i,x) = x_i$. Hence the
  multiplexer computes the index function with zero magic overhead.
\end{construction}

\begin{lemma}[Index function one-way / $\Rparpub$ lower bound]
  \label{lem:index-rac}
  \uses{def:index, def:smp}
  \notready
  $\Rparpub_\epsilon(\Index_n) \ge (1-h(\epsilon))\,n$.
\end{lemma}
% TODO: cited; follows from Nayak's optimal random-access-code lower bound [30].
% The R||^pub complexity is at least the one-way (Bob -> Alice) randomized
% complexity, which is >= (1-h(eps)) n since Bob's message defines a random
% access code; not reproved in the paper.

\begin{construction}[Controlled multiplexer recursion]
  \label{constr:controlled-multiplexer}
  \uses{def:multiplexer, def:toffoli}
  Define the \emph{controlled multiplexer}
  $\mathrm{cMultiplex}_n = \ketbra{0}{0}\otimes I + \ketbra{1}{1}\otimes
  \Multiplex_n$, acting on $2 + \lceil\log_2 n\rceil + n$ qubits; a controlled
  multiplexer implements an ordinary one by fixing the control to $\ket{1}$. The
  $2^{k+1}$-qubit controlled multiplexer is built from two $2^{k}$-qubit
  controlled multiplexers, with control $C$, index register
  $I_1,\dots,I_{k+1}$, array $X_1,\dots,X_{2^{k+1}}$, target $T$, and ancillas
  $A_1,A_2$:
  \begin{enumerate}
    \item apply $X$ to $I_1$; apply a Toffoli controlled on $C, I_1$ with target
      $A_1$; apply $X$ to $I_1$;
    \item apply a Toffoli controlled on $C, I_1$ with target $A_2$;
    \item run the $2^k$-qubit controlled multiplexer with control $A_1$, index
      $I_2,\dots,I_{k+1}$, first half $X_1,\dots,X_{2^k}$, target $T$;
    \item run the $2^k$-qubit controlled multiplexer with control $A_2$, index
      $I_2,\dots,I_{k+1}$, second half $X_{2^k+1},\dots,X_{2^{k+1}}$, target $T$;
    \item uncompute $A_1, A_2$.
  \end{enumerate}
  The ancillas $A_1, A_2$ record whether the active half is the left or right
  half (determined by the top index bit $I_1$ and control $C$); at most one of the
  two half-multiplexers fires, cleanly implementing the $2^{k+1}$-qubit controlled
  multiplexer.
\end{construction}

\begin{lemma}[Lemma 14: magic bounds for $\Multiplex_n$]
  \label{lem:multiplex-bound}
  \uses{def:multiplexer, def:magic-operator, def:binary-entropy}
  We have
  \[
    \Mmix_\epsilon(\Multiplex_n) \ge \frac{1}{4c_M}\big((1-h(\epsilon))\,n - 1
      \big),
  \]
  and the upper bound $\Mu_{\epsilon=0}(\Multiplex_n) = O(n)$.
\end{lemma}
\begin{proof}
  \uses{lem:unitary-computes-f, constr:multiplexer-index, thm:rpar-mixed,
    lem:index-rac, constr:controlled-multiplexer, def:binary-entropy}
  \emph{Lower bound.} By Construction~\ref{constr:multiplexer-index},
  $\Multiplex_n$ computes $\Index_n$ with no magic overhead, so by
  Lemma~\ref{lem:unitary-computes-f}, $\Mmix_\epsilon(\Multiplex_n) \ge
  \Mmix_\epsilon(\Index_n)$. By Theorem~\ref{thm:rpar-mixed},
  $\Mmix_\epsilon(\Index_n) \ge \tfrac{1}{4c_M}(\Rparpub_\epsilon(\Index_n) - 1)$,
  and by Lemma~\ref{lem:index-rac} $\Rparpub_\epsilon(\Index_n) \ge
  (1-h(\epsilon))n$, giving the stated bound.

  \emph{Upper bound.} Use Construction~\ref{constr:controlled-multiplexer}. Let
  $g(k)$ be the magic count (Toffoli gates of weight $3$) of the $2^k$-qubit
  controlled multiplexer. The recursion uses two copies of the $2^k$-qubit
  controlled multiplexer plus a constant number of Toffolis (steps 1, 2 and the
  uncomputation in step 5), so $g(k+1) \le 2g(k) + 4$, with base case $g(1) =
  O(1)$. Unrolling, $g(k) \le 2^k g_0 + 4(2^k - 1) = O(2^k)$. Since the
  $n$-array multiplexer corresponds to $k = \log_2 n$, the magic count is
  $O(2^{\log_2 n}) = O(n)$, and an ordinary multiplexer is a controlled one with
  control fixed to $\ket{1}$. Hence $\Mu_{\epsilon=0}(\Multiplex_n) = O(n)$.
\end{proof}

\chapter{Communication upper bounds from magic depth}

\section{The private simultaneous message passing model and NLQC}

\begin{theorem}[Theorem 15: low $T$-depth unitaries have efficient NLQC]
  \label{thm:speelman-nlqc}
  \uses{def:tdepth, def:nlqc}
  \notready
  Let $U_{AB}$ be a Clifford$+T$ unitary of $T$-depth at most $d$ acting on $n$
  qubits. Then $U_{AB}$ can be implemented as an NLQC using at most $O((68n)^d)$
  bits of communication and at most $O((68n)^d)$ shared EPR pairs.
\end{theorem}
% TODO: cited from Speelman [7] ("Instantaneous non-local computation of low
% T-depth quantum circuits"); the paper does not reprove it.

\begin{definition}[Garden-hose model and complexity $\GH(f)$]
  \label{def:garden-hose}
  \uses{def:teleportation}
  In the \emph{garden-hose model}, Alice and Bob share a fence with a number of
  pipes; Alice holds $x\in\{0,1\}^n$, Bob holds $y\in\{0,1\}^n$, Alice has a tap.
  Each player connects pipe-ends (and Alice connects the tap) by hoses, as a
  function of their input. Water from the tap flows along the resulting path and
  spills on one side; the side encodes the value of a Boolean function $f(x,y)$
  (spilling on Alice's side $=0$, Bob's side $=1$). The \emph{garden-hose
  complexity} $\GH(f)$ is the minimal number of pipes in a garden-hose protocol
  computing $f$. In the quantum picture, pipes are shared EPR pairs and hose
  connections are Bell-basis measurements, so the path describes concatenated
  teleportations (Definition~\ref{def:teleportation}).
\end{definition}

\begin{lemma}[Lemma 16: instantaneous garden-hose teleportation]
  \label{lem:gh-instantaneous}
  \uses{def:garden-hose, def:teleportation}
  \notready
  Let $f$ have garden-hose complexity $\GH(f)$, and suppose Alice initially holds
  $P^{f(x,y)}\ket{\psi}$ with $x$ known to Alice and $y$ known to Bob. Then:
  \begin{enumerate}
    \item there is an instantaneous protocol (no communication) using $2\GH(f)$
      EPR pairs, after which Alice holds $X^{g(\hat x)} Y^{h(\hat x)}\ket{\psi}$,
      where $\hat x$ consists of $x$ together with the $2\GH(f)$ bits describing
      Alice's and Bob's measurement outcomes;
    \item the garden-hose complexities of $g$ and $h$ are at most linear in
      $\GH(f)$:
      \[
        \GH(g) \le 4\GH(f) + 1, \qquad \GH(h) \le 11\GH(f) + 2.
      \]
  \end{enumerate}
\end{lemma}
% TODO: cited from Speelman [7]; not reproved in the paper.

\begin{lemma}[Lemma 17: garden-hose complexity of XORs]
  \label{lem:gh-xor}
  \uses{def:garden-hose}
  \notready
  Let $f_1,\dots,f_m$ be Boolean functions and $c\in\{0,1\}$. Then for $f =
  f_1\oplus\dots\oplus f_m\oplus c$ we have $\GH(f) \le 4\sum_{i=1}^m \GH(f_i) +
  1$.
\end{lemma}
% TODO: cited from Speelman [7]; not reproved in the paper.

\section{Transforming $\Qpars$ protocols into $\PSMs$ protocols}

\begin{theorem}[Theorem 18: $\Qpars \to \PSMs$ for low-$T$-depth referees]
  \label{thm:q-to-psm}
  \uses{def:smp, def:psm, def:tdepth}
  Consider an $\epsilon$-correct $\Qpars$ protocol for $f$ using $m$ qubits of
  message, in which the referee applies a $T$-depth-$d$ unitary to the received
  messages together with at most $a$ qubits of ancilla and then measures the
  first qubit. Then there is a $\PSMs_{\epsilon,\delta=2\epsilon}$ protocol for $f$
  using $O((68(m+a))^d)$ qubits of communication and entanglement.
\end{theorem}
\begin{proof}
  \uses{thm:speelman-nlqc, lem:gh-instantaneous, lem:gh-xor, def:psm,
    def:teleportation, def:tdepth, def:trace-norm}
  Suppose Alice and Bob have executed their $\Qpars$ actions and hold message
  systems $M_A, M_B$. Bob teleports $M_B$ to Alice (Definition~\ref{def:teleportation})
  without revealing the Pauli corrections. Alice attempts to execute the referee's
  unitary $U_{M_A M_B E}$ ($E$ the referee's advice) on her copy, but only knows
  Bob's state up to Pauli corrections; she tracks these through the Clifford$+T$
  decomposition of $U$.

  \emph{Clifford layers.} The teleportation produces Pauli corrections on the
  inputs; conjugating through a Clifford layer yields new Pauli corrections that
  are known to Bob (he knows the circuit).

  \emph{$T$ layers.} From the relations $TX = PXT$ and $TZ = ZT$
  (Definition~\ref{def:magic-gate}), Pauli $X$ corrections commute through a $T$
  gate at the cost of a conditional phase-gate $P$ correction. These conditional
  $P$ corrections are removed using garden-hose gadgets via
  Lemma~\ref{lem:gh-instantaneous}: the function deciding whether a $P$ correction
  is needed is initially fully known to Bob, hence has garden-hose complexity $1$;
  by Lemma~\ref{lem:gh-instantaneous} the resulting corrections after the first
  Clifford$+T$ layer remain of constant garden-hose complexity. For the next
  layer, the corrections conjugate through the Clifford to corrections whose
  presence is an XOR of a subset of previous corrections; by
  Lemma~\ref{lem:gh-xor} these have garden-hose complexity $O(m+a)$, and so on,
  the gadgets growing in cost layer by layer.

  \emph{Accounting.} The total entanglement cost of the gadgets matches the cost
  of implementing $U$ as an NLQC, namely $O((68(m+a))^d)$ by
  Theorem~\ref{thm:speelman-nlqc} with $n = m+a$. After the final Clifford$+T$
  layer Alice applies a last Clifford layer and measures a single qubit, with
  outcome bit $s$. Alice sends $s$ together with all her Bell-basis measurement
  outcomes to the referee; Bob sends all his Bell-basis measurement outcomes. From
  the Bell-basis outcomes the referee determines whether there was a Pauli $X$
  correction on the measured qubit and hence learns the corrected outcome,
  recovering $f(x,y)$ with error $\epsilon$. All messages are classical, so this
  is a $\Rpars$ (indeed $\PSMs$) protocol of cost $O((68(m+a))^d)$.

  \emph{Security ($\delta = 2\epsilon$).} Let $\vec r = (r_1,\dots,r_k)$ be the
  Bell-basis outcomes (uniformly random in $\{0,1\}^{|r|}$) and $s$ Alice's single
  output bit. The message has the form
  \[
    \rho_M(x,y) = \frac{1}{2^{|r|}}\sum_r X^{p(r)}\sigma(x,y)X^{p(r)}\otimes
      \ketbra{r}{r}, \quad
    \sigma(x,y) = \alpha(x,y)\ketbra{f}{f} + (1-\alpha(x,y))\ketbra{f\oplus1}{f\oplus1},
  \]
  where $p(r)$ is a parity deciding the $X$ correction on the measured qubit and
  $\alpha(x,y)\ge 1-\epsilon$. Define the simulator
  \[
    \Sim_M(f) = \frac{1}{2^{|r|}}\sum_r X^{p(r)}\ketbra{f}{f}X^{p(r)}\otimes
      \ketbra{r}{r}.
  \]
  Then
  \[
    \onenorm{\rho_M(x,y) - \Sim_M(f)}
      = \frac{1}{2^{|r|}}\sum_r \onenorm{\sigma(x,y) - \ketbra{f}{f}}
      = \frac{1}{2^{|r|}}\sum_r 2|1-\alpha(x,y)| \le 2\epsilon ,
  \]
  using unitary invariance of the trace norm and
  $\onenorm{(\alpha-1)\ketbra{f}{f} + (1-\alpha)\ketbra{f\oplus1}{f\oplus1}} =
  2|1-\alpha|$. Hence the protocol is $\delta = 2\epsilon$ secure.
\end{proof}

\begin{proposition}[Remark 19: simulation works for relational problems]
  \label{prop:relational-sim}
  \uses{thm:q-to-psm}
  \notready
  The simulation of $\Qpars$ by $\Rpars$ in Theorem~\ref{thm:q-to-psm} also
  applies to relational problems, dropping the privacy condition: the only change
  is that Alice sends the referee the measurement outcomes of a subset of qubits
  rather than a single qubit.
\end{proposition}
% TODO: stated as a remark without separate proof; same construction as
% Theorem 18 with a multi-qubit final measurement.

\begin{proposition}[Remark 20: $T$-depth lower bound from a $\PSMs$ vs $\Qpars$ gap]
  \label{prop:tdepth-lb}
  \uses{thm:q-to-psm, def:psm, def:smp, def:tdepth}
  \notready
  Let $\Lambda_f$ denote the measurement applied by the referee in a $\Qpars$
  protocol for $f$. Then Theorem~\ref{thm:q-to-psm} yields the $T$-depth lower
  bound
  \[
    T\text{-}\mathrm{depth}(\Lambda_f) = \Omega\!\left(
      \frac{\PSMs(f)}{\Qpars(f)}\right).
  \]
\end{proposition}
% TODO: stated as a remark; obtained by contrapositive from Theorem 18
% (if the referee had small T-depth, PSM*(f) would be polynomial in Q||*(f)).

\section{Separating $\Rpars$ and $\Rtwo$}

\begin{definition}[Definition 21: the Forrelation problem]
  \label{def:forrelation}
  For $x\in\{-1,1\}^n$ with $n$ a power of $2$, write $x = (x_1, x_2)$ (first and
  second halves) and define the \emph{forrelation}
  \[
    \forr(x) := \frac{1}{n}\,\bra{x_1} H^{\otimes n}\ket{x_2}.
  \]
  In the Forrelation communication problem Alice gets $x\in\{-1,1\}^n$, Bob gets
  $y\in\{-1,1\}^n$, and for a constant $\alpha>0$ the goal is to output
  \[
    f(x,y) = \begin{cases} -1 & \text{if } \forr(x\cdot y) \ge \alpha, \\
      +1 & \text{if } \forr(x\cdot y) \le \alpha/2, \end{cases}
  \]
  where $x\cdot y$ is the entrywise product.
\end{definition}

\begin{definition}[Definition 22: the ABCD problem]
  \label{def:abcd}
  Alice gets $A,C\in\mathrm{SU}(n)$ and Bob gets $B,D\in\mathrm{SU}(n)$, and the
  goal is to output
  \[
    f(x,y) = \begin{cases} -1 & \text{if } \Tr(ABCD) \ge 0.9\,n, \\
      +1 & \text{if } \Tr(ABCD) \le 0.1\,n. \end{cases}
  \]
\end{definition}

\begin{lemma}[Forrelation $\Rtwo$ lower bound]
  \label{lem:forr-r-lb}
  \uses{def:forrelation, def:two-way}
  \notready
  For $\alpha = \Theta(1/\log n)$, the Forrelation problem requires
  $\Rtwo$-cost $\tilde\Omega(n^{1/4})$.
\end{lemma}
% TODO: cited from Girish--May--Orshansky--Waddell [32] (building on
% Girish--Raz--Tal [14]); not reproved.

\begin{construction}[Constant-$T$-depth $\Qpars$ protocol for Forrelation]
  \label{constr:forr-protocol}
  \uses{def:forrelation, def:smp, def:tdepth}
  \notready
  There is a $\Qpars$ protocol of cost $O(\log n)$ for the Forrelation problem in
  which the referee's actions can be implemented in $T$-depth $2$.
\end{construction}
% TODO: cited Q||* upper bound from [32]; the constant-T-depth implementation of
% the referee is asserted there. Not reproved here.

\begin{lemma}[ABCD $\Rtwo$ lower bound]
  \label{lem:abcd-r-lb}
  \uses{def:abcd, def:two-way}
  \notready
  The ABCD problem requires $\Rtwo$-cost $\Omega(\sqrt{n})$.
\end{lemma}
% TODO: cited from Arunachalam--Girish--Lifshitz [15]; not reproved.

\begin{construction}[Constant-$T$-depth $\Qpars$ protocol for ABCD]
  \label{constr:abcd-protocol}
  \uses{def:abcd, def:smp, def:tdepth, def:toffoli}
  Alice and Bob share $\log n + 1$ EPR pairs. Alice applies
  $\bigl(\begin{smallmatrix} A & 0 \\ 0 & C \end{smallmatrix}\bigr)$ to her halves
  and Bob applies
  $\bigl(\begin{smallmatrix} B^\dagger & 0 \\ 0 & D^\dagger \end{smallmatrix}\bigr)$
  to his, and they send all qubits to the referee. The referee does a CNOT on the
  first two qubits, applies a controlled-swap between the last two register sets
  controlled on the first register, and measures the first qubit in the Hadamard
  basis. The referee outputs $1$ with probability $\ge 0.95$ if $\Tr(ABCD)\ge0.9n$
  and $\le 0.55$ if $\Tr(ABCD)\le0.1n$, solving the ABCD problem with cost
  $O(\log n)$. The referee's actions have $T$-depth $1$: the first CNOT is
  Clifford; before each controlled-swap the control qubit is copied onto
  $\log n$ ancillas by CNOTs (Clifford) so all CSWAPs run in parallel; each CSWAP
  is implemented with one Toffoli (Definition~\ref{def:toffoli}); and a Toffoli is
  implementable in $T$-depth $1$.
\end{construction}

\begin{theorem}[Exponential separation between $\Rpars$ and $\Rtwo$]
  \label{thm:separation-rpar-r}
  \uses{thm:q-to-psm, def:smp, def:two-way}
  There exist $n$-bit partial Boolean functions whose $\Rpars$ complexity is
  $\mathrm{polylog}(n)$ while their $\Rtwo$ (two-way randomized) complexity is
  $n^{\Omega(1)}$; in particular $\Rpars$ and $\Rtwo$ are exponentially
  separated.
\end{theorem}
\begin{proof}
  \uses{thm:q-to-psm, constr:abcd-protocol, lem:abcd-r-lb,
    constr:forr-protocol, lem:forr-r-lb}
  Take either the ABCD problem or the Forrelation problem
  (Definitions~\ref{def:abcd}, \ref{def:forrelation}). Each has a $\Qpars$
  protocol of cost $O(\log n)$ whose referee acts in constant $T$-depth
  (Constructions~\ref{constr:abcd-protocol} and \ref{constr:forr-protocol},
  with $T$-depth $1$ and $2$ respectively). Applying Theorem~\ref{thm:q-to-psm}
  with $m, a = O(\log n)$ and $d = O(1)$ converts each into an $\Rpars$ (in fact
  $\PSMs$) protocol of cost $(O(\log n))^{O(1)} = \mathrm{polylog}(n)$. On the
  other hand the $\Rtwo$ complexity is $\Omega(\sqrt n)$ for ABCD
  (Lemma~\ref{lem:abcd-r-lb}) and $\tilde\Omega(n^{1/4})$ for Forrelation
  (Lemma~\ref{lem:forr-r-lb}), both $n^{\Omega(1)}$. Hence $\Rpars =
  \mathrm{polylog}(n)$ while $\Rtwo = n^{\Omega(1)}$, an exponential separation.
\end{proof}
